\section{Introduction and main results}\label{sec:intro}

An on/off process alternates on periods and off periods with independent lengths. Let
the on periods have the law of $X_\plus$ and the off periods the law of $X_\minus$, two laws
on $\{1,2,\dots\}$ with regularly varying tails $T_\pm(n)=\Prob(X_\pm\ge n)$ of index
$\alpha\in(0,1)$, so that the periods have infinite mean, and let $A_N$ be the number of
time slots spent in the on state among the first $N$. Then $A_N/N$ does not concentrate. If
$T_\minus(n)/T_\plus(n)\to b\in(0,\infty)$, its limit law is the generalized arcsine law with
density
\begin{equation}\label{eq:lamperti-intro}
 \La(x)=\frac{\sin\pi\alpha}{\pi}\,
 \frac{b\,x^{\alpha-1}(1-x)^{\alpha-1}}{b^2x^{2\alpha}+2b\cos(\pi\alpha)\,x^\alpha(1-x)^\alpha+(1-x)^{2\alpha}},
 \qquad 0<x<1,
\end{equation}
which goes back to Lamperti~\cite{lamperti1958} and is called the Lamperti distribution in
the physics literature~\cite{margolinbarkai2006,zdk2015}. This paper is about the local form
of this law, at single points $k$, and about what happens at the ends of the range, under an
initial period of a different law, under aging, and when one tail is heavier than the other.
The same process is a lattice L\'evy walk in the ballistic regime, which moves in runs of
alternating direction and sits at $2A_N-N$ after $N$ steps, so $P_N(k)=\Prob(A_N=k)$ is also
the law of the walk on its natural lattice. We keep the vocabulary of the walk in the body
of the paper. A period is a run, an on slot is a plus step, and the two ends of the range
are the plus end $k=N$ and the minus end $k=0$. The companion paper~\cite{paperF} treats the
finite-mean case $1<\alpha<2$, where $A_N-pN$ is stable at the scale $N^{1/\alpha}$, and the two
papers share the exact identities of Section~\ref{sec:setting}.

\subsection{Main results}\label{sec:intro-results}

\emph{The local Lamperti law.} A renewal law with a tail $T$ of index $\alpha\in(0,1)$
satisfies the strong renewal theorem if its renewal mass function $u$ satisfies
$u(n)\sim\frac{\sin\pi\alpha}\pi\frac1{nT(n)}$. For $\alpha>1/2$ this always holds under
aperiodicity, by Garsia and Lamperti~\cite{garsialamperti1963} and
Erickson~\cite{erickson1970}. For $\alpha\le1/2$ it can fail, and Caravenna and
Doney~\cite{caravennadoney2019} found the exact condition. Our first theorem says that the
local Lamperti law is equivalent to the strong renewal theorem for both run laws.

\begin{introtheorem}[the local Lamperti law]\label{thm:local}
Assume that $T_\plus$ and $T_\minus$ are regularly varying with index $-\alpha$,
$0<\alpha<1$, that $T_\minus(n)/T_\plus(n)\to b\in(0,\infty)$, and that the run laws
are aperiodic. Let the walk start with a fresh run of either sign. Then the
following are equivalent.
\begin{enumerate}[(i)]
\item $NP_N(\lfloor x_0N\rfloor)\to\La(x_0)$ for one $x_0\in(0,1)$.
\item $NP_N(k)=\La(k/N)+o(1)$ uniformly over $k$ with $k/N$ in each compact
subset of $(0,1)$.
\item Both run laws satisfy the strong renewal theorem.
\end{enumerate}
\end{introtheorem}

With the theorem of Caravenna and Doney this gives a complete answer. For $\alpha>1/2$ the
local law holds under regular variation and aperiodicity alone. For $\alpha\le1/2$ it holds
if and only if each run law satisfies their condition, which we recall in
Lemma~\ref{lem:cd}, and Doney's local condition $f_\pm(n)=O(T_\pm(n)/n)$ suffices. A single
interior point already forces the strong renewal theorem for both laws. Without it the lower
half of the local law still holds, $\liminf_N(NP_N(k)-\La(k/N))\ge0$ uniformly on compact
sets, and the failure shows up as an excess at every interior point
(Corollary~\ref{cor:lower}). The proof does not need a bivariate renewal theorem. An exact
identity writes $P_N(k)$ through the convolution powers of the two run laws, sorted by the
number $m$ of completed runs. For $m$ of order $1/T_\plus(N)$ the local limit theorem of
Gnedenko gives a Riemann sum for an explicit double integral. The key estimate,
Lemma~\ref{lem:key}, bounds the probability that the $(m+1)$-th run covers a level $k$ by
$CT(k)$ uniformly in $m$, under regular variation alone. It decouples the two signs for small
$m$, and the contribution of small $m$ of each sign is then exactly the quantity that
Caravenna and Doney show to be negligible if and only if the strong renewal theorem holds.
We identify the limit with $\La$ through its Stieltjes transform
(Proposition~\ref{prop:lambda}), which gives a self-contained derivation of Lamperti's density.

\emph{The ends.} At distance $j$ from the plus end the law is the tail of a run times a
renewal mass function, $P_N(N-j)=T_\plus(N)u_\minus(j)+o(T_\plus(N))$, under regular variation
of $T_\plus$ alone. Under Doney's condition this boundary layer matches the bulk uniformly up
to distance $o(N)$ from the end (Theorem~\ref{thm:ends}), and under the strong renewal theorem
alone the upper half of the matching holds for $\alpha\ne1/2$ (Lemma~\ref{lem:E}). Since
$NT_\plus(N)\to\infty$, the end values are much larger than the bulk values, so the global mode
lies at bounded distance from an end for every $\alpha\in(0,1)$ (Corollary~\ref{cor:mode}).
For $1<\alpha<2$ the same comparison has a golden ratio threshold~\cite{paperF}, and here it
has none.

\emph{The initial run.} Wang, Schulz, Deng and Barkai~\cite{wsdb2018} left ``the effect of
initial preparation on the rare fluctuations'' for future work~\cite[Section~VIII]{wsdb2018}.
For the occupation count with $1<\alpha<2$ this is answered in~\cite{paperF}, and
Theorem~\ref{thm:initial} answers it for $0<\alpha<1$. Let the first run have sign $\sgn$ and
length $d$ with probability $\nu_\sgn(d)$. In the finite-mean case the bulk needs no condition
on the initial run on any macroscopic interior. Here it does. A first run that ends just
before the $k$-th plus step leaves a remainder that sits in its own boundary layer, where the
probabilities are of order $T(N)$ and not $1/N$. Theorem~\ref{thm:initial} gives the exact
defect in the bulk, the necessary condition $\nu_\pm(n)=o(1/(nT_\plus(n)))$, the sufficient
condition $n\nu_\pm(n)\to0$ under Doney's condition, and an exact criterion. It also shows by
example that an initial run which is finite almost surely can change the local law at
interior points, and that the necessary condition is not sufficient. At the ends the
criterion becomes a window condition.

\emph{Initial runs of length of order $N$, and aging.} If the law of the initial run depends
on $N$ and puts mass of order $1$ on lengths of order $N$, the local law changes to a mixture
of rescaled Lamperti densities. Theorem~\ref{thm:mixture} proves this under the strong renewal
theorem for both laws. The near-diagonal terms, the initial runs that end within $o(N)$ steps
of the point, are controlled by a counting bound (Lemma~\ref{lem:count}) and not by the
end-layer theorem, which is why Doney's condition is not needed. The main application is
aging. Let the process start fresh at time $0$ and count the plus steps among the steps
$t_a+1,\dots,t_a+N$, with $t_a/N\to\theta$. The run in progress at time $t_a$ is an initial run
of length of order $N$, with the marked Dynkin--Lamperti law of Lemma~\ref{lem:DL}, and
Theorem~\ref{thm:aging} gives the aged law at every point. In the bulk it is the mixture
$L_\theta$ of~\eqref{eq:L-theta}, which is the aging generalized arcsine law of Akimoto, Sera,
Yamato and Yano~\cite{asyy2020}, there obtained in distribution by a scaling argument. At the
endpoints there are atoms of order $1$, with the closed form~\eqref{eq:aged-constants}. In the
end layers the shape $u_\minus(j)$ of the fresh start survives and its height is multiplied by
an explicit constant $c_\plus(\theta)<1$, with $c_\plus(1)=1/(2\sqrt2)$ for $\alpha=1/2$ and
equal tails. The limits $\theta\to0$ and $N\to\infty$ commute in the bulk and not at the ends.

\emph{One tail heavier.} When $T_\minus=o(T_\plus)$ and $\alpha_\plus<1$, the minus side may have
an infinite mean with any index $\alpha_\minus\in(0,1)$ or a finite mean with
$\alpha_\minus\in(1,2)$. Then $A_N/N\to1$, and the object is the time $N-A_N$ on the lighter
side, which lives at the scale $D_N=\mu_\minus/T_\plus(N)$ when the lighter mean is finite and at
the scale $D_N$ given by $T_\minus(D_N)\sim T_\plus(N)$ when it is infinite. The limit in
distribution is a Mittag-Leffler law in the first case and the law of $E^{1/\alpha_\minus}Z_\minus$
in the second, with $E$ Mittag-Leffler and $Z_\minus$ stable, and it is known in distribution
through Darling and Kac~\cite{darlingkac1957}, C\'enac, Le Ny, de Loynes and
Offret~\cite[Section~4.3]{clnlo2019} and Yamato~\cite{yamato2026}. Theorem~\ref{thm:dominated}
gives the local law. It holds under regular variation and aperiodicity alone when the lighter
mean is finite or $\alpha_\minus>1/2$, and for $\alpha_\minus\le1/2$ exactly when the lighter
law satisfies the strong renewal theorem and the heavier law satisfies a weak condition
(W$_\plus$), which holds under its own strong renewal theorem. The strong renewal theorem of the
heavier law never enters. The end layer at the heavier end is $T_\plus(N)u_\minus(j)$ and joins
the bulk at the same height, and the layer at the lighter end is $T_\minus(N)u_\plus(i)$. The
proof needs one new estimate, the local form of the inverse renewal count of the heavier runs
(Lemma~\ref{lem:Q}), and a uniform one-big-jump estimate for sums of a growing number of
heavy-tailed summands (Lemma~\ref{lem:onejump}), which we prove by a Chernoff bound. Together
with~\cite{paperF}, this covers every pair of indices in $(0,1)\cup(1,2)$ for which
$T_\minus/T_\plus$ has a limit.

\emph{The shape.} For equal tails the density $L_{\alpha,1}$ is U-shaped for small $\alpha$ and
W-shaped for $\alpha$ close to $1$, with the transition at $\alpha_c=0.5946\ldots$, the root of
$\cos(\pi\alpha/2)=\alpha$. The threshold is in Baldassarri, Bouchaud, Dornic and
Godr\`eche~\cite{bbdg1999} and in Godr\`eche and Luck~\cite{godrecheluck2001}.
Theorem~\ref{thm:shape} classifies the critical points of $\La$ for every $\alpha$ and every
$b$. The density has $3$ critical points if and only if $|\log b|<\beta_c(\alpha)$, where
$\beta_c$ is given in closed form, vanishes for $\alpha\le\alpha_c$ and is infinite for
$\alpha\ge\alpha_1=0.7365\ldots$, the root of $\sin(\pi\alpha)=\alpha$, so that beyond $\alpha_1$
the W shape survives every asymmetry. The proof reduces the critical points to the level set
of one function of one variable whose monotonicity we determine. At $b=0$ the dominated regime
takes over, and its ratio density gains an interior maximum exactly at $\alpha_1$, so the two
shape thresholds meet. The discrete law has more extrema than the density. The end profiles
add local maxima of their own, and Theorem~\ref{thm:signs} gives the lower bound
$\sigma_L+\sigma_\plus+\sigma_\minus$ on the number of sign changes of the first differences,
with $\sigma_L\in\{1,3\}$ from Theorem~\ref{thm:shape} and $\sigma_\pm$ from the profiles.
Proposition~\ref{prop:count} shows that two difference statements, a difference form of
Doney's local bound in the layer and a second-difference local limit theorem in the bulk,
turn the bound into the exact count $M_\plus+M_\minus+m_L$ of modes for all large $N$, with
$m_L\in\{0,1\}$ the number of interior maxima of $\La$. We prove the pieces of the count that
do not need these statements, in particular a difference form of the strong renewal theorem
for log-convex laws (Lemma~\ref{lem:dsrt}), and we leave the two statements open. In every
computed case the bound is attained.

\subsection{Related work}\label{sec:intro-related}

\emph{Occupation times.} Lamperti~\cite{lamperti1958} proved an occupation time theorem for
a class of processes with a regeneration state, which allows two sojourn laws with different
weights, and the density $\La$ is attached to his name. Mitov and Yanev~\cite{mitovyanev2001}
study limit theorems for alternating renewal processes with infinite mean cycles, extending
results of Dynkin, Lamperti and Erickson. Tak\'acs~\cite{takacs1957} and
Kesten~\cite{kesten1962} proved the limit laws under finite means, and Starreveld, Bekker and
Mandjes~\cite{sbm2018} give the transform of the occupation time of an alternating renewal
process at an exponential time, a central limit theorem under finite second moments, and large
deviations under exponential moments. Godr\`eche and Luck~\cite{godrecheluck2001} give the
transform of the occupation time and its limit laws for one sojourn law, including the
threshold $\alpha_c$ for the shape. All of these are statements in distribution. We found no
local limit theorem for the occupation count of an alternating renewal process with infinite
means, and no treatment of the dominated regime at the local level.

\emph{Strong renewal theorems.} The strong renewal theorem for $\alpha>1/2$ is due to Garsia
and Lamperti~\cite{garsialamperti1963}, Williamson~\cite{williamson1968} and
Erickson~\cite{erickson1970}. Doney~\cite{doney1997} proved it for all $\alpha\in(0,1)$ under
the local condition $f(n)=O(T(n)/n)$, together with a local large deviation bound. Caravenna
and Doney~\cite{caravennadoney2019} gave the necessary and sufficient condition for
$\alpha\le1/2$, and their local large deviation bound for sums of a fixed number of variables
is the main input of our key estimate. Berger~\cite{berger2019ptrf} extends Doney's local
bound, and his version is the one we quote. Berger~\cite{berger2019} also proved multivariate
strong renewal theorems. Theorem~\ref{thm:local} shows that for the occupation count the
bivariate problem reduces exactly to the two univariate ones.

\emph{Physics.} In the physics literature the walk is the L\'evy walk in the ballistic
regime~\cite{zdk2015}, and the occupation count is the time average of a blinking signal with
power-law sojourn times~\cite{margolinbarkai2006}. Both sources use $L_{\alpha,1}$ as the limit
density. Wang, Schulz, Deng and Barkai~\cite[Section~VII.A]{wsdb2018} found an atom with a
boundary layer for the occupation time when $0<\alpha<1$, for one sojourn law and a start of
fair random sign, through a double Laplace transform, and Godr\`eche and
Luck~\cite[Section~7]{godrecheluck2001} have the same structure. Theorem~\ref{thm:ends} is a
lattice version of this layer at a fixed time, for two laws, with the uniform matching to the
bulk. Aging of renewal processes with power-law sojourns is a standard theme there, and
Akimoto, Sera, Yamato and Yano~\cite{asyy2020} proved the aged arcsine law for Brownian motion
and stated its generalization to two-state renewal processes with a scaling argument.
Akimoto~\cite{akimoto2023} computes moments of the aged occupation time. Theorem~\ref{thm:aging}
is the local form, with the atoms and the end layers.

\emph{The finite-mean case.} For $1<\alpha<2$ the occupation law has a stable bulk, big-jump
zones and boundary layers, and its initial-run dependence, its aging and its modes are
described in~\cite{paperF}. We use the exact identities of that paper and recall them with
proofs in Section~\ref{sec:identities}. The two papers together describe the local law of
$A_N$ on the square $(0,1)\cup(1,2)$ of index pairs, up to the lines $\alpha_\sgn=1$, which
Section~\ref{sec:open} leaves open.

\subsection{Structure of the paper}\label{sec:intro-structure}

Section~\ref{sec:setting} fixes notation and proves the exact identities.
Section~\ref{sec:tools} collects the cited theorems and proves the elementary bounds, the key
estimate, the one-big-jump estimate and the local inverse renewal count.
Section~\ref{sec:local} proves Theorem~\ref{thm:local}, Section~\ref{sec:ends} treats the ends
and the global mode, Section~\ref{sec:initial} the initial run, Section~\ref{sec:aging} the
mixture law and aging, Section~\ref{sec:dominated} the dominated regime, and
Section~\ref{sec:shape} the shape. Section~\ref{sec:open} lists the open problems and
Section~\ref{sec:verify} describes the computations. No proof uses a computed number.
